\documentclass[11pt]{article}
\usepackage[margin=0.85in]{geometry}
\usepackage[T1]{fontenc}
\usepackage{hyperref}
\begin{document}
\sloppy
\section*{Independent signed graph cycle review}
No blocking discrepancy found at the scalar-field cycle and anchor scope. Fresh isolated replay passes ten producer theorems, eight producer controls, eight reviewer controls and 26 separate axiom readbacks, with only propext, Classical.choice and Quot.sound. Three copied modules match producer bytes; all producer manifest entries match. The producer was not modified. Final main1,controls1,reviewer3,readback1,reviewer\_readback1 pass with unchanged bindings.\par\medskip
Source-first derivation matches native three\_leads.rs1597-1613 and1659-1712. After potential reduction, a same-root edge is (cu+cv)*root=a+b*k. Opposite signs yield a+b*k=0; native passes -a,-b to constraint, so its returned -(-a)/(-b) equals k. Equal signs yield root=a/(2cu)+b/(2cu)*k, requiring nonzero2 and nonzero sign. The native denominator branch 2 versus r-2 corresponds to2cu for signs plus/minus one. Existing fixed root x+y*k compared to a+b*k yields (x-a)+(y-b)*k=0 and extraction -(x-a)/(y-b) when the latter coefficient is nonzero. With zero coefficient, sound equations imply zero numerator, matching the assertion condition.\par\medskip
Reviewer controls over ZMod7 apply original theorems to opposite signs, negative equal signs and fixed-anchor differences. They reject using positive2 for the negative-sign anchor, check zero consistent/inconsistent constraints, and test the minus-sign denominator. A characteristic-two control shows an actual anchor failure if the nonzero2 premise is dropped: root1 satisfies the doubled zero equation but the inverse formula returns0. First reviewer attempt needed explicit field instances and typed arguments; second needed simplification before deciding noncomputable inverse expressions. Both failed receipts remain, and the corrected controls use kernel evaluation and algebraic rewrites.\par\medskip
Native bound nuance: Ring::new alone checks prime r>2 and does not impose a word-addition bound. The inspected specific pruned.rs caller at line814 asserts2<=r<=2\textasciicircum{}63. With canonical residues that implies a+b<=2r-2<2\textasciicircum{}64, so a blanket claim that this caller lacks an addition bound would be wrong. Other callers and canonicality/inversion correspondence must be audited separately. The current proof is scalar Field algebra; machine i8/u64 reflection, forest state and node coordinate interpretation are outside its scope.\par\medskip
Review sequencing and independence: the source-first PLAN was saved before reading producer Cycles.lean in this review, using prior knowledge of the graph algorithm. The inputs were sealed private artifacts, not published/precommitted review inputs; the Coordinator records that sequencing deviation and is publishing their archive under T-1225. This report does not claim a precommitted plan or fully blind project history. The reviewer authored imported GraphAlgebra earlier, while the new cycle module was authored separately; the independent finding concerns these new cycle/sign/constraint statements. Imported binaries retain their recorded trust boundary. No scientific execution, repository edits, proof promotion or compiled-equivalence claim occurred.\par\medskip
\paragraph{Artifact references} Cycles.lean; Controls.lean; ReviewerControls.lean; status.json. The editable diagram is flow.svg. Source-first derivation and review deviations are in PLAN.md and REPORT.md.
\end{document}
